\subsection{Likelihood tilting and compact-local quadratic risk}
\label{subsec:v34-tilting-moments}

We give the likelihood-tilting and moment arguments underlying
Lemma~\ref{lem:v22-contiguity} and the estimator assertions of
Theorem~\ref{thm:v22-vector-boundary-gaussian}.  This also specifies how the
unit-mean normalization in Lemma~\ref{lem:v33-finite-likelihood} is used
under local alternatives.  The argument concerns the moving-support model
\eqref{eq:v22-thinned-family}; physical acquisition is handled by the
separate stopped-transfer results.

Three distinct estimates enter the conclusion.  The density expansion
\eqref{eq:v35-one-mark-expansion} gives the alternative mean
\eqref{eq:v35-alternative-mean}; centering under that same alternative
then gives \eqref{eq:v35-centered-fourth}.  These are direct moment
estimates for the original product law.  The normalized likelihoods and
the LAN expansion instead give the weak shifted limit
\eqref{eq:v34-direct-tilting}.  Combining that weak limit with the
original-law fourth-moment bound yields
\eqref{eq:v34-quadratic-risk}.  Thus the argument for unbounded quadratic
loss does not rely on total-variation comparison alone.  The further
passage from finite parameter sets to compact experiments uses the
uniform modulus in Theorem~\ref{thm:v32-compact-vector}.

\begin{proposition}[Likelihood tilting and uniform quadratic risk]
\label{prop:v34-tilting-moments}
Under the hypotheses of Theorem~\ref{thm:v22-vector-boundary-gaussian},
let \(h_n\to h\) be a bounded local sequence.  Under the reference-dominated
common-collar experiment, set
\[
 L_n=\frac{d\overline P_{n,h_n}^{\otimes n}}
          {d\overline P_{n,0}^{\otimes n}}.
\]
The likelihoods \(L_n\) are uniformly integrable under their reference
measures.  The reference and alternative sequences are mutually contiguous,
and, for every bounded continuous \(\varphi\),
\begin{equation}\label{eq:v34-direct-tilting}
 \mathbb E_{\overline P_{n,h_n}^{\otimes n}}\varphi(\Delta_n)
 \longrightarrow
 \mathbb E\!\left[
  e^{h^tZ-h^t\mathcal J_\Sigma h/2}\varphi(Z)\right],
 \qquad Z\sim N(0,\mathcal J_\Sigma).
\end{equation}
The limit is the expectation under
\(N(\mathcal J_\Sigma h,\mathcal J_\Sigma)\), including singular
information.  The same central-sequence limit holds in the original
experiment, with observations outside \(C_n\) contributing zero.

For every compact \(K\subset\mathbb R^r\), the original product experiment
satisfies
\begin{equation}\label{eq:v34-fourth-moment}
 \sup_{n\ge n_0(K)}\sup_{h\in K}
        \mathbb E_{P_{n,h}^{\otimes n}}\|\Delta_n\|^4<\infty.
\end{equation}
Write \(J=\mathcal J_\Sigma\), \(\mathsf H=\operatorname{Ran}J\), and
\(\Pi=\Pi_{\mathsf H}\).  If \(W\) is positive semidefinite on
\(\mathsf H\), then
\begin{equation}\label{eq:v34-quadratic-risk}
 \sup_{h\in K}\left|
 \mathbb E_{P_{n,h}^{\otimes n}}
  \bigl[(J^+\Delta_n-\Pi h)^tW(J^+\Delta_n-\Pi h)\bigr]
                 -\operatorname{tr}(WJ^+)\right|\longrightarrow0.
\end{equation}
Here \(J^+\) acts only on the identifiable space; no risk assertion for an
unidentifiable coordinate is made.
\end{proposition}

\begin{proof}
The uniform LAN expansion gives joint convergence under the reference law,
\[
 (\Delta_n,L_n)\Rightarrow
 (Z,L),\qquad L=e^{h^tZ-h^tJh/2}>0.
\]
Both \(\mathbb E L_n\) and \(\mathbb E L\) equal one.  For fixed
\(A>0\), convergence of the bounded continuous truncation gives
\[
 \mathbb E(L_n-A)_+=1-\mathbb E(L_n\wedge A)
        \longrightarrow1-\mathbb E(L\wedge A)=\mathbb E(L-A)_+.
\]
Since \(L_n\mathbf1_{\{L_n>2A\}}\le2(L_n-A)_+\), this proves uniform
integrability, without any assumption on higher likelihood moments.
Apply joint weak convergence to
\((L_n\wedge A)\varphi(\Delta_n)\), then let \(A\to\infty\).
The tail bound just proved yields \eqref{eq:v34-direct-tilting}.
Completing the Gaussian exponential, or computing its characteristic
function on \(\operatorname{Ran}J\), identifies the displayed shifted
normal law.  This calculation does not invert \(J\) on its kernel.

For completeness, let \(P_n\) and \(Q_n=L_nP_n\) denote the censored
reference and alternative product measures.  For every event \(A_n\),
\[
 Q_n(A_n)\le A P_n(A_n)
           +\mathbb E_{P_n}[L_n\mathbf1_{\{L_n>A\}}].
\]
Uniform integrability proves contiguity of \(Q_n\) with respect to \(P_n\).
In the reverse direction, for every \(c>0\),
\[
 P_n(A_n)\le P_n\{L_n\le c\}+c^{-1}Q_n(A_n).
\]
First take a continuity value \(c\) for the law of \(L\), and then send
\(c\downarrow0\).  Positivity of \(L\) proves reverse contiguity.
Subsequence extraction proves the assertion for every bounded sequence,
not only a convergent one.  Embed the original and censored observations in
their common space with the extra symbol \(\partial\).  Their laws differ
only on an event of probability at most \(C_Knp_nq_n^2=o(1)\), uniformly
on \(K\).  This transfers the contiguity assertions to the original
experiments.  The central sequence is unchanged by censoring, since every
censored observation already contributes zero to its definition.

We next prove the moment estimate directly under the original alternatives.
On the common collar \(C_n\), joint smoothness and positivity give the
\emph{one-successful-observation} density bound
\begin{equation}\label{eq:v34-one-mark-ratio}
 \left|\frac{f_{\delta_nh}}{f_0}-1\right|
       \le C_K\delta_n(1+w_0^{-1}),\qquad
 \sup_{h\in K}\left\|\frac{f_{\delta_nh}}{f_0}-1\right\|_{L^\infty(C_n)}
       =O_K(\delta_n/q_n)=o(1).
\end{equation}
This is not a bound tending to one for the product likelihood.
The collar estimates and \eqref{eq:v34-one-mark-ratio} imply
\begin{align}
 \left\|\int_{C_n}\mathsf S f_{\delta_nh}\,dx\right\|
       &\le Cq_n+C_K\delta_n\log(1/q_n),
                                      \label{eq:v34-mean-budget}\\
 \int_{C_n}\|\mathsf S\|^2 f_{\delta_nh}\,dx
       &\le C_K\log(1/q_n),\qquad
 \int_{C_n}\|\mathsf S\|^4 f_{\delta_nh}\,dx\le C_Kq_n^{-2}.
                                      \label{eq:v34-moment-budget}
\end{align}
For the first inequality, the reference truncated mean is \(O(q_n)\).
In normal coordinates \(s=w_0\), the remaining integrand is bounded by
\(C_K\delta_n(s+2+s^{-1})\); integration from \(q_n\) to a fixed collar
width gives the logarithm.  The complement of the collar is uniformly
regular.  The other inequalities follow from the uniformly bounded density
ratio and Lemma~\ref{lem:v22-vector-collar}.

Put \(\xi_{n,i}=\delta_nY_i\mathsf S(X_i)\mathbf1_{C_n}(X_i)\).
These vectors are independent under every fixed alternative.  Their sum is
\(\Delta_n\), and \eqref{eq:v34-mean-budget} gives
\[
 \sup_{h\in K}\|\mathbb E_{n,h}\Delta_n\|
 \le Cnp_n\delta_nq_n+C_Knp_n\delta_n^2\log(1/q_n)=O_K(1).
\]
For centered independent vectors in fixed dimension, expansion of the fourth
power, coordinate by coordinate, yields
\begin{align*}
 \mathbb E_{n,h}\|\Delta_n-\mathbb E_{n,h}\Delta_n\|^4
 &\le C_r\left\{
       \left(\sum_{i=1}^n\mathbb E_{n,h}\|\xi_{n,i}\|^2\right)^2
              +\sum_{i=1}^n\mathbb E_{n,h}\|\xi_{n,i}\|^4\right\}\\
 &\le C_K\left\{
     [np_n\delta_n^2\log(1/q_n)]^2+np_n\delta_n^4q_n^{-2}\right\}
 =O_K(1).
\end{align*}
Indeed, with \(q_n=\delta_n[\log(1/\delta_n)]^{1/4}\), the first mean
term \(np_n\delta_nq_n\) and the final fourth-moment term tend to zero,
whereas \(np_n\delta_n^2\log(1/q_n)\to1\).
This proves \eqref{eq:v34-fourth-moment}; each excluded observation, rather
than the entire product sample on a censoring event, contributes zero.

The tilted central-sequence limit and \eqref{eq:v34-fourth-moment} imply
convergence of quadratic risks along every convergent local sequence.
The Gaussian error \(J^+Z-\Pi h\), for
\(Z\sim N(Jh,J)\), has covariance \(J^+\) and mean zero.
Its \(W\)-risk is therefore \(\operatorname{tr}(WJ^+)\).
If \eqref{eq:v34-quadratic-risk} were not uniform on \(K\), a violating
sequence would have a convergent subsequence in \(K\), contradicting the
just-proved sequential limit.  This proves the compact-uniform assertion.
\end{proof}
